% Source: https://artofproblemsolving.com/community/c3h2482930_judgefan99s_problem_set
% Stable identity: aops-topic-2482930-q02
Since $8^2+15^2=64+225=289=17^2$, the triangle is right-angled, with hypotenuse $17$ (by the converse of Pythagoras' theorem).

\textbf{Circumradius.} The hypotenuse of a right triangle is a diameter of its circumcircle (the right angle is an angle in a semicircle), so $R=\dfrac{17}{2}$.

\textbf{Inradius.} The area is $K=\dfrac12\cdot8\cdot15=60$ and the semiperimeter is $s=\dfrac{8+15+17}{2}=20$. From $K=rs$ we get $r=3$.

Hence
\[R+r=\dfrac{17}{2}+3=\boxed{\dfrac{23}{2}}.\]

\begin{asy}
usepackage("amsmath");
size(6.5cm);
defaultpen(fontsize(10pt));
pair C=(0,0), B=(15,0), A=(0,8);
pair Oc=(A+B)/2, I=(3,3);
draw(circle(Oc,8.5),gray);
draw(circle(I,3),blue);
draw(A--B--C--cycle,linewidth(0.8));
draw((0,0.8)--(0.8,0.8)--(0.8,0));
draw(Oc--B,red); draw(I--(3,0),blue+dashed);
dot(Oc,red); dot(I,blue);
label("$A$",A,NW); label("$B$",B,E); label("$C$",C,SW);
label("$8$",(A+C)/2,W); label("$15$",(10,0),S); label("$17$",A+0.28*(B-A),NE);
label("$O$",Oc,N,red); label("$R=\dfrac{17}{2}$",(Oc+B)/2,N,red);
label("$I$",I,N,blue); label("$r=3$",(3,1.5),E,blue);
\end{asy}
